%-----------------------------------------------------------------------------------------------------------------------------------------------%
%	The MIT License (MIT)
%
%	Copyright (c) 2015 Jan Küster
%
%	Permission is hereby granted, free of charge, to any person obtaining a copy
%	of this software and associated documentation files (the "Software"), to deal
%	in the Software without restriction, including without limitation the rights
%	to use, copy, modify, merge, publish, distribute, sublicense, and/or sell
%	copies of the Software, and to permit persons to whom the Software is
%	furnished to do so, subject to the following conditions:
%	
%	THE SOFTWARE IS PROVIDED "AS IS", WITHOUT WARRANTY OF ANY KIND, EXPRESS OR
%	IMPLIED, INCLUDING BUT NOT LIMITED TO THE WARRANTIES OF MERCHANTABILITY,
%	FITNESS FOR A PARTICULAR PURPOSE AND NONINFRINGEMENT. IN NO EVENT SHALL THE
%	AUTHORS OR COPYRIGHT HOLDERS BE LIABLE FOR ANY CLAIM, DAMAGES OR OTHER
%	LIABILITY, WHETHER IN AN ACTION OF CONTRACT, TORT OR OTHERWISE, ARISING FROM,
%	OUT OF OR IN CONNECTION WITH THE SOFTWARE OR THE USE OR OTHER DEALINGS IN
%	THE SOFTWARE.
%	
%
%-----------------------------------------------------------------------------------------------------------------------------------------------%


%============================================================================%
%
%	DOCUMENT DEFINITION
%
%============================================================================%

%we use article class because we want to fully customize the page and dont use a cv template
\documentclass[10pt, Letter]{article}	




%----------------------------------------------------------------------------------------
%	ENCODING
%----------------------------------------------------------------------------------------

%we use utf8 since we want to build from any machine
\usepackage[utf8]{inputenc}		

%----------------------------------------------------------------------------------------
%	LOGIC
%----------------------------------------------------------------------------------------

% provides \isempty test
\usepackage{xifthen}

%----------------------------------------------------------------------------------------
%	FONT
%----------------------------------------------------------------------------------------

% some tex-live fonts - choose your own

%\usepackage[defaultsans]{droidsans}
%\usepackage[default]{comfortaa}
%\usepackage{cmbright}
\usepackage[default]{raleway}
%\usepackage{fetamont}
%\usepackage[default]{gillius}
%\usepackage[light,math]{iwona}
%\usepackage[thin]{roboto} 

% set font default
\renewcommand*\familydefault{\sfdefault} 	
\usepackage[T1]{fontenc}

% more font size definitions
\usepackage{moresize}		

% href package
\usepackage{hyperref}
%----------------------------------------------------------------------------------------
%	PAGE LAYOUT  DEFINITIONS
%----------------------------------------------------------------------------------------

%debug page outer frames
%\usepackage{showframe}			


%define page styles using geometry
\usepackage[a4paper]{geometry}		

% for example, change the margins to 2 inches all round
\geometry{top=1.25cm, bottom=-.6cm, left=1.5cm, right=1.5cm} 	

%use customized header
\usepackage{fancyhdr}				
\pagestyle{fancy}

%less space between header and content
\setlength{\headheight}{-5pt}		


%customize entries left, center and right
%\lhead{}
%\chead{}
%\rhead{}


%indentation is zero
\setlength{\parindent}{0mm}

%----------------------------------------------------------------------------------------
%	TABLE /ARRAY DEFINITIONS
%---------------------------------------------------------------------------------------- 

%for layouting tables
\usepackage{multicol}			
\usepackage{multirow}

%extended aligning of tabular cells
\usepackage{array}

\newcolumntype{x}[1]{%
>{\raggedleft\hspace{0pt}}p{#1}}%


%----------------------------------------------------------------------------------------
%	GRAPHICS DEFINITIONS
%---------------------------------------------------------------------------------------- 

%for header image
\usepackage{graphicx}

%for floating figures
\usepackage{wrapfig}
\usepackage{float}
%\floatstyle{boxed} 
%\restylefloat{figure}

%for drawing graphics		
\usepackage{tikz}				
\usetikzlibrary{shapes, backgrounds,mindmap, trees}


%----------------------------------------------------------------------------------------
%	Color DEFINITIONS
%---------------------------------------------------------------------------------------- 

\usepackage{color}

%accent color
\definecolor{sectcol}{RGB}{0,96,163}

%dark background color
\definecolor{bgcol}{RGB}{110,110,110}

%light background / accent color
\definecolor{softcol}{RGB}{225,225,225}


%============================================================================%
%
%
%	DEFINITIONS
%
%
%============================================================================%

%----------------------------------------------------------------------------------------
% 	HEADER
%----------------------------------------------------------------------------------------

% remove top header line
\renewcommand{\headrulewidth}{0pt} 

%remove botttom header line
\renewcommand{\footrulewidth}{0pt}	  	

%remove pagenum
\renewcommand{\thepage}{}	

%remove section num		
\renewcommand{\thesection}{}			

%----------------------------------------------------------------------------------------
% 	ARROW GRAPHICS in Tikz
%----------------------------------------------------------------------------------------

% a six pointed arrow poiting to the left
\newcommand{\tzlarrow}{(0,0) -- (0.2,0) -- (0.3,0.2) -- (0.2,0.4) -- (0,0.4) -- (0.1,0.2) -- cycle;}	

% include the left arrow into a tikz picture
% param1: fill color
%
\newcommand{\larrow}[1]
{\begin{tikzpicture}[scale=0.58]
	 \filldraw[fill=#1!100,draw=#1!100!black]  \tzlarrow
 \end{tikzpicture}
}

% a six pointed arrow poiting to the right
\newcommand{\tzrarrow}{ (0,0.2) -- (0.1,0) -- (0.3,0) -- (0.2,0.2) -- (0.3,0.4) -- (0.1,0.4) -- cycle;}

% include the right arrow into a tikz picture
% param1: fill color
%
\newcommand{\rarrow}
{\begin{tikzpicture}[scale=0.7]
	\filldraw[fill=sectcol!100,draw=sectcol!100!black] \tzrarrow
 \end{tikzpicture}
}



%----------------------------------------------------------------------------------------
%	custom sections
%----------------------------------------------------------------------------------------

% create a coloured box with arrow and title as cv section headline
% param 1: section title
%
\newcommand{\cvsection}[1]
{
	\begin{center}
		\large\textcolor{sectcol}{\textbf{#1}}
	\end{center}
}

%create a coloured arrow with title as cv meta section section
% param 1: meta section title
%
\newcommand{\metasection}[2]
{
%\begin{tabular*}{1\textwidth}{r r}
\footnotesize{#2} \hspace*{\fill} \footnotesize{#1}\\[1pt]
%\end{tabular*}
}

%----------------------------------------------------------------------------------------
%	 CV EVENT
%----------------------------------------------------------------------------------------

% creates a stretched box as cv entry headline followed by some paragraphs about 
% the work you did
% param 1:	event time i.e. 2014 or 2011-2014 etc.
% param 2:	event name (what did you do?)
% param 3:	institution (where did you work / study)
% param 4:	list of paragraphs
%
\newcommand{\cvevent}[4]
{

\begin{tabular*}{1\textwidth}{p{13.6cm}  x{3.9cm}}
	\textbf{#2} - \textcolor{bgcol}{#3} &   \vspace{2.5pt}\textcolor{sectcol}{#1}
\end{tabular*}

\vspace{-8pt}
\textcolor{softcol}{\hrule}
\vspace{6pt}

	\foreach \desc in {#4}{
		$\cdot$ \desc\\[3pt]
	}
	
\vspace{3pt}
}

% creates a stretched box as 
\newcommand{\cveventmeta}[2]
{
	\mbox{\mystrut \hspace{87pt}\textit{#1}}\\
	#2
}

%----------------------------------------------------------------------------------------
% CUSTOM STRUT FOR EMPTY BOXES
%----------------------------------------- -----------------------------------------------
\newcommand{\mystrut}{\rule[-.3\baselineskip]{0pt}{\baselineskip}}

%----------------------------------------------------------------------------------------
% CUSTOM LOREM IPSUM
%----------------------------------------------------------------------------------------
\newcommand{\lorem}
{Lorem ipsum dolor sit amet, consectetur adipiscing elit. Donec a diam lectus.}



%============================================================================%
%
%
%
%	DOCUMENT CONTENT
%
%
%
%============================================================================%
\begin{document}


%use our custom fancy header definitions
\pagestyle{fancy}	





%----------------------------------------------------------------------------------------
%	HEADER IMAGE
%----------------------------------------------------------------------------------------

%\begin{figure}[H]
%\begin{flushright}
%	\includegraphics[trim= 320 130 460 210,clip,width=0.2\linewidth]{myfoto.jpg}	%trimming relative to image size!
%\end{flushright}
%\end{figure}


%---------------------------------------------------------------------------------------
%	TITLE HEADLINE
%----------------------------------------------------------------------------------------
\vspace{-8pt}
\begin{center}
	\huge \textsc{Asif Al Noor} \\[2pt]
	\small RF Systems Design | System Design (Avionics Hardware) | Assembly, Integration \& Test (AIT)
\end{center}



\vspace{6pt}


%---------------------------------------------------------------------------------------
%	META SECTION
%----------------------------------------------------------------------------------------
\vspace{-10pt}
\metasection{Toronto, Canada (Canadian Citizen)}{roonlafisa@gmail.com | +1-905-783-3906 | \href{https://roonlafisa.github.io}{https://roonlafisa.github.io}}
\vspace{-6pt}
\textcolor{softcol}{\hrule}


\normalsize

%---------------------------------------------------------------------------------------
%	SUMMARY (optional)
%----------------------------------------------------------------------------------------

\cvsection{Summary}
\vspace{-5pt}
I have around 8 years of RF system design experience for IoT, satellites, and rockets. More recently, I led the design and development of a full-stack (design to assembly) avionics system for an suborbital-scale rocket. 
\vspace{-2pt}
\cvsection{Skills}
\vspace{-15pt}
\begin{multicols}{2}
\textbf{System design}: End-to-end RF subsystem, Full-stack hardware subsystem, PCB design, JIRA. \\
\textbf{RF design}: RF simulations, RF system architecture, Link budgets, Ground station, Antenna design, Filter design. \\
\textbf{RF hardware}: RFD900x, Iridium, Globalstar, Quasonix, Safran, CCSDS Receivers, Ublox GPS, JAVAD GPS, LimeSDR, uAvionics Transponder. \\
\textbf{RF software}: CST, HFSS, ADS, COMSOL, MATLAB tools.

\columnbreak

\textbf{Standards}: IPC-610, IPC-620, AS50881G MIL-810, MIL-883, MIL-3899, MIL-83513, MIL-24308. \\
\textbf{Electronics}: Design w/ESP32, STM32; implementing i2C, CAN, RS-232, TTL. \\
\textbf{Programming}: Python, C, MATLAB, Cube IDE, Linux. \\
\textbf{PCB design}: LT Spice, Altium, Circuit Studio, KiCAD. \\
\textbf{CAD prototyping}: Solidworks, Solidworks electrical. \\
\textbf{Lab}: Oscilloscope, Signal generator, Electrical load, Network Analyzer, Logic analyzer, SMD soldering. 
\end{multicols}

%============================================================================%
%
%	CV SECTIONS AND EVENTS (MAIN CONTENT)
%
%============================================================================%

%---------------------------------------------------------------------------------------
%	EXPERIENCE
%----------------------------------------------------------------------------------------

\cvsection{Professional History}
\vspace{-5pt}
\cvevent{Sept 2021 - Feb 2023}{Senior RF Specialist}{\textbf{SpaceRyde}, Toronto ON}{
	{Designed, implemented and manufactured the \textbf{avionics system of a rocket} including flight computers, GNC sensosrs, power management, data networking, RF communications on subsystem and component levels.},
   % {Designed, built and commissioned the \textbf{attitude control system} of a rocket launch platform from stratospheric balloon including its through motor, ESC, propellers.},
	{Architected and commissioned the \textbf{end-to-end RF system} to receive rocket telemetry using legacy COTS design. Scope  of responsibilities also included the \textbf{ground station} implementation and operation.},
%	{Developed \textbf{RF link budgeting} a dynamic software from scratch and verified it with real-time test data.},
    {Designed, assembled and tested \textbf{PCBs} incorporating MCU, serial devices, and other critical components.},
%	{CAD-modeled and assembled avionics flight components, wires, and connectors inside flight enclosures.},
	{Spearheaded the development of \textbf{test procedures and campaigns} to validate the avionics components.},
	{Produced \textbf{systems requirements} and derived requirements for developing avionics components.},
        {Interfaced with regulatory authority (ISED) regarding \textbf{radio licensing} for rocket telemetry spectrum licensing.}
	%{Maintained an excellent relationship with tier-1 suppliers in negotiating and procuring space-grade hardware.}
}

%
\cvevent{April 2020 - Aug 2021}{Intermediate RF Engineer}{\textbf{SMT Research}, Vancouver BC}{
	{Designed and implemented a battery-less \textbf{RF power harvesting} wireless system for structural monitoring.},
	{Pioneered custom \textbf{RFID passive sensor} tags for remote data logging and pinpointing hidden sensor locations.},
%    {Designed, simulated and implemented \textbf{patch antenna at 915MHz} for remote data logging with RFID tag.},
	{Installed sensors and DAQs in active construction sites and provided technical support to installation technicians.}
}

%
\cvevent{ July 2017 - May 2019}{RF Engineering Consultant}{\textbf{Direct Kinetic Solutions}, Texas (Remote)}{
	{Produced custom \textbf{Cubesat EPS and RF system} concepts for clients and DOD grant applications.}
}


%
\cvevent{ July 2017 - May 2019}{Lead RF Engineer}{\textbf{Helios Wire}, Vancouver BC}{
	{Designed, simulated and successfully deployed \textbf{C, X and S-Band antennas} for CubeSat TT\&C communication.},
	{Conceptualized and modeled an iot gateway to interface between Cubesat in space and IoT tags in the ground.}
}
%---------------------------------------------------------------------------------------
%	EDUCATION SECTION
%--------------------------------------------------------------------------------------
\vspace{-20pt}
\cvsection{Educational History}
\vspace{-5pt}
\cvevent{Sept 2014 - Oct 2016}{MASc in EE}{University of British Columbia, Kelowna BC}{
	{Master Thesis: ''A broadband fixed-beam leaky-wave antenna based on transformation electromagnetics.''}
}

%
\cvevent{ Jan 2010 - Oct 2013}{BSc in EEE}{Islamic University of Technology, Bangladesh}{
	{Activities: Led the university robotics team during the NASA Lunabotics Annual Competition 2013.}
}
%---------------------------------------------------------------------------------------
%	EDUCATION SECTION
%--------------------------------------------------------------------------------------
\vspace{-20pt}
\cvsection{Ongoing Personal Projects}
\vspace{-15pt}
\begin{multicols}{2}
\textbf{Embedded Hardware and Firmware}: Developing a development kit using STM32F407 with modular shields. \\

\columnbreak

\textbf{Computer vision}: Implementing a gesture detection-based game on Tx2i Jetson Nano for my toddler. \\
 
\end{multicols}


%============================================================================%
%
%
%
%	DOCUMENT END
%
%
%
%============================================================================%
\end{document}